% ---- Acknowledgements + Appendices + References ----------------------

\section*{Acknowledgements}

We thank J.-W. Chen, K.~Cranmer, W.~Detmold, R.~Melko, D.~Murphy, A.~Pochinsky, and B.~Trippe for helpful discussions.
%
This work is supported in part by the U.S.~Department of Energy, Office of Science, Office of Nuclear Physics under grant Contract Number DE-SC0011090. This research used resources of the Argonne Leadership Computing Facility, which is a DOE Office of Science User Facility supported under Contract DE-AC02-06CH11357, under the ALCF Aurora Early Science Program. Some work was undertaken at the Kavli Institute for Theoretical Physics, supported by the National Science Foundation under Grant No.~NSF PHY-1748958.
PES is supported by the National Science Foundation under CAREER Award 1841699, GK is supported by the U.S.~Department of Energy under the SciDAC4 award DE-SC0018121, and PES and MSA are supported in part by NSERC and the Perimeter Institute for Theoretical Physics. Research at the Perimeter Institute is supported by the Government of Canada through the Department of Innovation, Science and Economic Development and by the Province of Ontario through the Ministry of Research and Innovation.


\appendix

\section{Acceptance rate estimator for autocorrelation}
\label{app:bounds}

Here it is shown that the standard estimator for the autocorrelation of an observable $\mathcal{O}$ converges in the limit of infinite path length to  $p_{\tau \mathrm{rej}}$, as claimed in Section~\ref{subsec:autocorr}. The standard estimator is defined by:
%
\begin{align}
& \lim_{N\rightarrow\infty} \expt{}\sq{\widehat{\rho(\tau)/\rho(0)}_O} \\
=\,&
    \lim_{N\rightarrow\infty} \expt{} \sq{\frac{\frac{1}{N-\tau} \sum_{i=0}^{N-\tau-1} (\obs_i - \bar{\obs})(\obs_{i+\tau}-\bar{\obs})}{\frac{1}{N} \sum_{i=0}^{N-1} \paren{\obs_i - \bar{\obs}}^2}} \\
=\,& \lim_{N\rightarrow\infty} \expt{} \sq{\frac{(\obs_i - \bar{\obs})(\obs_{i+\tau}-\bar{\obs})}{\frac{1}{N} \sum_{i=0}^{N-1} \paren{\obs_i - \bar{\obs}}^2}},
\end{align}
%
where the final equality is true assuming the Markov chain is initialized with a sample from the stationary distribution $p(\phi)$ (i.e., it is assumed that enough prior iterations were discarded that the chain is thermalized). The expectation value can then be split by cases and, conditioning on the fixed accept/reject pattern, the expectation values can be computed by identifying the distributions of observables $\mathcal{O}_i$ and $\mathcal{O}_{i+\tau}$:
%
% \begin{widetext} removed (single-column normalized layout)
\begin{align}
\lim_{N\rightarrow\infty} \expt{}\sq{\widehat{\rho(\tau)/\rho(0)}_O}
&= \lim_{N\rightarrow \infty} p_{\tau\text{rej}} \expt{} \sq{\frac{(\obs_i - \bar{\obs})(\obs_{i+\tau}-\bar{\obs})}{\frac{1}{N} \sum_{i=0}^{N-1} \paren{\obs_i - \bar{\obs}}^2}
    \Bigg| \text{all proposals } i+1, \dots, i+\tau \text{ rejected}} + \nonumber\\
    &\qquad\qquad (1-p_{\tau\text{rej}}) \expt{} \sq{\frac{(\obs_i - \bar{\obs})(\obs_{i+\tau}-\bar{\obs})}{\frac{1}{N} \sum_{i=0}^{N-1} \paren{\obs_i - \bar{\obs}}^2}
    \Bigg| \text{some proposal } i+1, \dots, i+\tau \text{ accepted}} \\
&= p_{\tau\text{rej}} \lim_{N\rightarrow\infty} \expt{\phi\sim p} \sq{\frac{(\obs[\phi] - \bar{\obs})^2}{\frac{1}{N} \sum_{j \notin [i,i+\tau]} (\obs_j - \bar{\obs})^2 + O(\tau/N)}} + \nonumber\\
    &\qquad\qquad (1-p_{\tau\text{rej}}) \lim_{N\rightarrow\infty} \expt{\phi\sim p} \expt{\phi' \sim \tilde{p}} \sq{ \frac{(\obs[\phi] - \bar{\obs})(\obs[\phi'] - \bar{\obs})}{\frac{1}{N} \sum_{j \notin [i,i+\tau]} (\obs_j - \bar{\obs})^2 + O(\tau/N)}} \\
&= p_{\tau\text{rej}} \frac{\rho(0)}{\rho(0)} + (1-p_{\tau\text{rej}}) \frac{\cancel{\expt{\phi\sim p} \sq{\obs[\phi] - \bar{\obs}}} \expt{\phi'\sim \tilde{p}} \sq{\obs[\phi'] - \bar{\obs}}}{\rho(0)} = p_{\tau\text{rej}}.
\end{align}
% \end{widetext}
%
The limit $N \rightarrow \infty$ is used to drop biases arising from conditioning on behavior within the region $[i,i+\tau]$.


\section{Mean absolute error loss}
\label{app:mae-loss}
The mean absolute error (MAE) loss optimized before training some models is defined by:
%
\begin{equation}
L_\text{MAE}(\netp) := \int \prod_j d\phi_j \, \netp(\phi) \left| \log{\netp(\phi)} - \log{p(\phi)} \right|.
\end{equation}
%
It is bounded below by the KL divergence $D_{KL}(\netp || p)$ and has global minima exactly where the KL loss does, when $\netp = p$.

In practice, the loss is stochastically estimated by drawing batches of $M$ samples from the model $\curly{\phi^{(i)} \sim \netp}$ and computing the sample mean:
%
\begin{equation}
\begin{aligned}
\widehat{L_\text{MAE}(\netp)} &= \frac{1}{M} \sum_{i=1}^M \left| \log{\netp(\phi^{(i)})} - \log{p(\phi^{(i)})} \right| \\
&= \frac{1}{M} \sum_{i=1}^M \left| \log{\netp(\phi^{(i)})} + S(\phi^{(i)}) + \log{Z} \right|.
\end{aligned}
\end{equation}
To employ this loss, the partition function must either be estimated ahead of time, or initialized as a trainable parameter. In this study, a multistage method~\cite{ValleauFreeEnergy72} was used to estimate and fix the partition function value used while optimizing $L_\text{MAE}$.

\setlength{\parskip}{2pt}
This loss is appealing due to the point-by-point potential driving the distribution towards the correct one. Any errors in computing $\log{Z}$, however, result in a minimum at which the model distribution $\netp(\phi)$ does not necessarily agree with the desired distribution $p(\phi)$. This loss was therefore only used prior to training with the shifted KL loss.

% ---- References (from the original main.bbl, unchanged) ---------------
\begin{thebibliography}{58}
\expandafter\ifx\csname natexlab\endcsname\relax\def\natexlab#1{#1}\fi
\expandafter\ifx\csname bibnamefont\endcsname\relax
  \def\bibnamefont#1{#1}\fi
\expandafter\ifx\csname bibfnamefont\endcsname\relax
  \def\bibfnamefont#1{#1}\fi
\expandafter\ifx\csname citenamefont\endcsname\relax
  \def\citenamefont#1{#1}\fi
\expandafter\ifx\csname url\endcsname\relax
  \def\url#1{\texttt{#1}}\fi
\expandafter\ifx\csname urlprefix\endcsname\relax\def\urlprefix{URL }\fi
\providecommand{\bibinfo}[2]{#2}
\providecommand{\eprint}[2][]{\url{#2}}

\bibitem[{\citenamefont{Wolff}(1990)}]{Wolff:1989wq}
\bibinfo{author}{\bibfnamefont{U.}~\bibnamefont{Wolff}},
  \bibinfo{journal}{Nucl. Phys. Proc. Suppl.} \textbf{\bibinfo{volume}{17}},
  \bibinfo{pages}{93} (\bibinfo{year}{1990}).

\bibitem[{\citenamefont{Del~Debbio et~al.}(2004)\citenamefont{Del~Debbio,
  Manca, and Vicari}}]{DelDebbio:2004xh}
\bibinfo{author}{\bibfnamefont{L.}~\bibnamefont{Del~Debbio}},
  \bibinfo{author}{\bibfnamefont{G.~M.} \bibnamefont{Manca}}, \bibnamefont{and}
  \bibinfo{author}{\bibfnamefont{E.}~\bibnamefont{Vicari}},
  \bibinfo{journal}{Phys. Lett.} \textbf{\bibinfo{volume}{B594}},
  \bibinfo{pages}{315} (\bibinfo{year}{2004}), \eprint{hep-lat/0403001}.

\bibitem[{\citenamefont{Meyer et~al.}(2007)\citenamefont{Meyer, Simma, Sommer,
  Della~Morte, Witzel, and Wolff}}]{Meyer:2006ty}
\bibinfo{author}{\bibfnamefont{H.~B.} \bibnamefont{Meyer}},
  \bibinfo{author}{\bibfnamefont{H.}~\bibnamefont{Simma}},
  \bibinfo{author}{\bibfnamefont{R.}~\bibnamefont{Sommer}},
  \bibinfo{author}{\bibfnamefont{M.}~\bibnamefont{Della~Morte}},
  \bibinfo{author}{\bibfnamefont{O.}~\bibnamefont{Witzel}}, \bibnamefont{and}
  \bibinfo{author}{\bibfnamefont{U.}~\bibnamefont{Wolff}},
  \bibinfo{journal}{Comput. Phys. Commun.} \textbf{\bibinfo{volume}{176}},
  \bibinfo{pages}{91} (\bibinfo{year}{2007}), \eprint{hep-lat/0606004}.

\bibitem[{\citenamefont{Kandel et~al.}(1988)\citenamefont{Kandel, Domany, Ron,
  Brandt, and Loh}}]{Kandel:1988zz}
\bibinfo{author}{\bibfnamefont{D.}~\bibnamefont{Kandel}},
  \bibinfo{author}{\bibfnamefont{E.}~\bibnamefont{Domany}},
  \bibinfo{author}{\bibfnamefont{D.}~\bibnamefont{Ron}},
  \bibinfo{author}{\bibfnamefont{A.}~\bibnamefont{Brandt}}, \bibnamefont{and}
  \bibinfo{author}{\bibfnamefont{E.}~\bibnamefont{Loh}},
  \bibinfo{journal}{Phys. Rev. Lett.} \textbf{\bibinfo{volume}{60}},
  \bibinfo{pages}{1591} (\bibinfo{year}{1988}).

\bibitem[{\citenamefont{Bonati and D'Elia}(2018)}]{Bonati:2017woi}
\bibinfo{author}{\bibfnamefont{C.}~\bibnamefont{Bonati}} \bibnamefont{and}
  \bibinfo{author}{\bibfnamefont{M.}~\bibnamefont{D'Elia}},
  \bibinfo{journal}{Phys. Rev.} \textbf{\bibinfo{volume}{E98}},
  \bibinfo{pages}{013308} (\bibinfo{year}{2018}), \eprint{1709.10034}.

\bibitem[{\citenamefont{Hasenbusch and Schaefer}(2018)}]{PhysRevD.98.054502}
\bibinfo{author}{\bibfnamefont{M.}~\bibnamefont{Hasenbusch}} \bibnamefont{and}
  \bibinfo{author}{\bibfnamefont{S.}~\bibnamefont{Schaefer}},
  \bibinfo{journal}{Phys. Rev. D} \textbf{\bibinfo{volume}{98}},
  \bibinfo{pages}{054502} (\bibinfo{year}{2018}),
  \urlprefix\url{https://link.aps.org/doi/10.1103/PhysRevD.98.054502}.

\bibitem[{\citenamefont{Swendsen and Wang}(1987)}]{SwendsenWang87}
\bibinfo{author}{\bibfnamefont{R.~H.} \bibnamefont{Swendsen}} \bibnamefont{and}
  \bibinfo{author}{\bibfnamefont{J.-S.} \bibnamefont{Wang}},
  \bibinfo{journal}{Phys. Rev. Lett.} \textbf{\bibinfo{volume}{58}},
  \bibinfo{pages}{86} (\bibinfo{year}{1987}),
  \urlprefix\url{https://link.aps.org/doi/10.1103/PhysRevLett.58.86}.

\bibitem[{\citenamefont{Wolff}(1989)}]{Wolff1989}
\bibinfo{author}{\bibfnamefont{U.}~\bibnamefont{Wolff}},
  \bibinfo{journal}{Phys. Rev. Lett.} \textbf{\bibinfo{volume}{62}},
  \bibinfo{pages}{361} (\bibinfo{year}{1989}),
  \urlprefix\url{https://link.aps.org/doi/10.1103/PhysRevLett.62.361}.

\bibitem[{\citenamefont{Prokof'ev and Svistunov}(2001)}]{ProkofevSvistunov01}
\bibinfo{author}{\bibfnamefont{N.}~\bibnamefont{Prokof'ev}} \bibnamefont{and}
  \bibinfo{author}{\bibfnamefont{B.}~\bibnamefont{Svistunov}},
  \bibinfo{journal}{Phys. Rev. Lett.} \textbf{\bibinfo{volume}{87}},
  \bibinfo{pages}{160601} (\bibinfo{year}{2001}),
  \urlprefix\url{https://link.aps.org/doi/10.1103/PhysRevLett.87.160601}.

\bibitem[{\citenamefont{{Kawashima} and {Harada}}(2004)}]{KawashimaHarada04}
\bibinfo{author}{\bibfnamefont{N.}~\bibnamefont{{Kawashima}}} \bibnamefont{and}
  \bibinfo{author}{\bibfnamefont{K.}~\bibnamefont{{Harada}}},
  \bibinfo{journal}{Journal of the Physical Society of Japan}
  \textbf{\bibinfo{volume}{73}}, \bibinfo{pages}{1379} (\bibinfo{year}{2004}),
  \eprint{cond-mat/0312675}.

\bibitem[{\citenamefont{Bietenholz et~al.}(1995)\citenamefont{Bietenholz,
  Pochinsky, and Wiese}}]{Bietenholz1995MeronCluster}
\bibinfo{author}{\bibfnamefont{W.}~\bibnamefont{Bietenholz}},
  \bibinfo{author}{\bibfnamefont{A.}~\bibnamefont{Pochinsky}},
  \bibnamefont{and} \bibinfo{author}{\bibfnamefont{U.~J.} \bibnamefont{Wiese}},
  \bibinfo{journal}{Phys. Rev. Lett.} \textbf{\bibinfo{volume}{75}},
  \bibinfo{pages}{4524} (\bibinfo{year}{1995}),
  \urlprefix\url{https://link.aps.org/doi/10.1103/PhysRevLett.75.4524}.

\bibitem[{\citenamefont{Ramos}(2012)}]{Ramos:2012bb}
\bibinfo{author}{\bibfnamefont{A.}~\bibnamefont{Ramos}}, \bibinfo{journal}{PoS}
  \textbf{\bibinfo{volume}{LATTICE2012}}, \bibinfo{pages}{193}
  (\bibinfo{year}{2012}), \eprint{1212.3800}.

\bibitem[{\citenamefont{Gambhir and Orginos}(2015)}]{Gambhir:2015rha}
\bibinfo{author}{\bibfnamefont{A.~S.} \bibnamefont{Gambhir}} \bibnamefont{and}
  \bibinfo{author}{\bibfnamefont{K.}~\bibnamefont{Orginos}},
  \bibinfo{journal}{PoS} \textbf{\bibinfo{volume}{LATTICE2014}},
  \bibinfo{pages}{043} (\bibinfo{year}{2015}), \eprint{1506.06118}.

\bibitem[{\citenamefont{Cossu et~al.}(2018)\citenamefont{Cossu, Boyle, Christ,
  Jung, Jüttner, and Sanfilippo}}]{Cossu:2017eys}
\bibinfo{author}{\bibfnamefont{G.}~\bibnamefont{Cossu}},
  \bibinfo{author}{\bibfnamefont{P.}~\bibnamefont{Boyle}},
  \bibinfo{author}{\bibfnamefont{N.}~\bibnamefont{Christ}},
  \bibinfo{author}{\bibfnamefont{C.}~\bibnamefont{Jung}},
  \bibinfo{author}{\bibfnamefont{A.}~\bibnamefont{Jüttner}}, \bibnamefont{and}
  \bibinfo{author}{\bibfnamefont{F.}~\bibnamefont{Sanfilippo}},
  \bibinfo{journal}{EPJ Web Conf.} \textbf{\bibinfo{volume}{175}},
  \bibinfo{pages}{02008} (\bibinfo{year}{2018}), \eprint{1710.07036}.

\bibitem[{\citenamefont{Jin and Osborn}(2019)}]{Jin2019QNHMC}
\bibinfo{author}{\bibfnamefont{X.-Y.} \bibnamefont{Jin}} \bibnamefont{and}
  \bibinfo{author}{\bibfnamefont{J.~C.} \bibnamefont{Osborn}}, in
  \emph{\bibinfo{booktitle}{Proceedings of Science}} (\bibinfo{year}{2019}),
  \eprint{1904.10039}.

\bibitem[{\citenamefont{Endres et~al.}(2015)\citenamefont{Endres, Brower,
  Detmold, Orginos, and Pochinsky}}]{Endres:2015yca}
\bibinfo{author}{\bibfnamefont{M.~G.} \bibnamefont{Endres}},
  \bibinfo{author}{\bibfnamefont{R.~C.} \bibnamefont{Brower}},
  \bibinfo{author}{\bibfnamefont{W.}~\bibnamefont{Detmold}},
  \bibinfo{author}{\bibfnamefont{K.}~\bibnamefont{Orginos}}, \bibnamefont{and}
  \bibinfo{author}{\bibfnamefont{A.~V.} \bibnamefont{Pochinsky}},
  \bibinfo{journal}{Phys. Rev.} \textbf{\bibinfo{volume}{D92}},
  \bibinfo{pages}{114516} (\bibinfo{year}{2015}), \eprint{1510.04675}.

\bibitem[{\citenamefont{Detmold and Endres}(2016)}]{Detmold:2016rnh}
\bibinfo{author}{\bibfnamefont{W.}~\bibnamefont{Detmold}} \bibnamefont{and}
  \bibinfo{author}{\bibfnamefont{M.~G.} \bibnamefont{Endres}},
  \bibinfo{journal}{Phys. Rev.} \textbf{\bibinfo{volume}{D94}},
  \bibinfo{pages}{114502} (\bibinfo{year}{2016}), \eprint{1605.09650}.

\bibitem[{\citenamefont{Detmold and Endres}(2018)}]{Detmold:2018zgk}
\bibinfo{author}{\bibfnamefont{W.}~\bibnamefont{Detmold}} \bibnamefont{and}
  \bibinfo{author}{\bibfnamefont{M.~G.} \bibnamefont{Endres}},
  \bibinfo{journal}{Phys. Rev.} \textbf{\bibinfo{volume}{D97}},
  \bibinfo{pages}{074507} (\bibinfo{year}{2018}), \eprint{1801.06132}.

\bibitem[{\citenamefont{Luscher and Schaefer}(2013)}]{Luscher:2012av}
\bibinfo{author}{\bibfnamefont{M.}~\bibnamefont{Luscher}} \bibnamefont{and}
  \bibinfo{author}{\bibfnamefont{S.}~\bibnamefont{Schaefer}},
  \bibinfo{journal}{Comput. Phys. Commun.} \textbf{\bibinfo{volume}{184}},
  \bibinfo{pages}{519} (\bibinfo{year}{2013}), \eprint{1206.2809}.

\bibitem[{\citenamefont{Mages et~al.}(2017)\citenamefont{Mages, Toth, Borsanyi,
  Fodor, Katz, and Szabo}}]{Mages:2015scv}
\bibinfo{author}{\bibfnamefont{S.}~\bibnamefont{Mages}},
  \bibinfo{author}{\bibfnamefont{B.~C.} \bibnamefont{Toth}},
  \bibinfo{author}{\bibfnamefont{S.}~\bibnamefont{Borsanyi}},
  \bibinfo{author}{\bibfnamefont{Z.}~\bibnamefont{Fodor}},
  \bibinfo{author}{\bibfnamefont{S.~D.} \bibnamefont{Katz}}, \bibnamefont{and}
  \bibinfo{author}{\bibfnamefont{K.~K.} \bibnamefont{Szabo}},
  \bibinfo{journal}{Phys. Rev.} \textbf{\bibinfo{volume}{D95}},
  \bibinfo{pages}{094512} (\bibinfo{year}{2017}), \eprint{1512.06804}.

\bibitem[{\citenamefont{Burnier et~al.}(2018)\citenamefont{Burnier, Florio,
  Kaczmarek, and Mazur}}]{Burnier:2017osu}
\bibinfo{author}{\bibfnamefont{Y.}~\bibnamefont{Burnier}},
  \bibinfo{author}{\bibfnamefont{A.}~\bibnamefont{Florio}},
  \bibinfo{author}{\bibfnamefont{O.}~\bibnamefont{Kaczmarek}},
  \bibnamefont{and} \bibinfo{author}{\bibfnamefont{L.}~\bibnamefont{Mazur}},
  \bibinfo{journal}{EPJ Web Conf.} \textbf{\bibinfo{volume}{175}},
  \bibinfo{pages}{07004} (\bibinfo{year}{2018}), \eprint{1710.06472}.

\bibitem[{\citenamefont{Laio et~al.}(2016)\citenamefont{Laio, Martinelli, and
  Sanfilippo}}]{Laio:2015era}
\bibinfo{author}{\bibfnamefont{A.}~\bibnamefont{Laio}},
  \bibinfo{author}{\bibfnamefont{G.}~\bibnamefont{Martinelli}},
  \bibnamefont{and}
  \bibinfo{author}{\bibfnamefont{F.}~\bibnamefont{Sanfilippo}},
  \bibinfo{journal}{JHEP} \textbf{\bibinfo{volume}{07}}, \bibinfo{pages}{089}
  (\bibinfo{year}{2016}), \eprint{1508.07270}.

\bibitem[{\citenamefont{Tanaka and Tomiya}(2017)}]{Tanaka:2017niz}
\bibinfo{author}{\bibfnamefont{A.}~\bibnamefont{Tanaka}} \bibnamefont{and}
  \bibinfo{author}{\bibfnamefont{A.}~\bibnamefont{Tomiya}}
  (\bibinfo{year}{2017}), \eprint{1712.03893}.

\bibitem[{\citenamefont{Shanahan et~al.}(2018)\citenamefont{Shanahan,
  Trewartha, and Detmold}}]{Shanahan:2018vcv}
\bibinfo{author}{\bibfnamefont{P.~E.} \bibnamefont{Shanahan}},
  \bibinfo{author}{\bibfnamefont{D.}~\bibnamefont{Trewartha}},
  \bibnamefont{and} \bibinfo{author}{\bibfnamefont{W.}~\bibnamefont{Detmold}},
  \bibinfo{journal}{Phys. Rev.} \textbf{\bibinfo{volume}{D97}},
  \bibinfo{pages}{094506} (\bibinfo{year}{2018}), \eprint{1801.05784}.

\bibitem[{\citenamefont{Huang and Wang}(2017)}]{Huang2017}
\bibinfo{author}{\bibfnamefont{L.}~\bibnamefont{Huang}} \bibnamefont{and}
  \bibinfo{author}{\bibfnamefont{L.}~\bibnamefont{Wang}},
  \bibinfo{journal}{Physical Review B} \textbf{\bibinfo{volume}{95}},
  \bibinfo{pages}{1} (\bibinfo{year}{2017}), ISSN \bibinfo{issn}{24699969}.

\bibitem[{\citenamefont{Liu et~al.}(2017)\citenamefont{Liu, Qi, Meng, and
  Fu}}]{Liu2017SLMC}
\bibinfo{author}{\bibfnamefont{J.}~\bibnamefont{Liu}},
  \bibinfo{author}{\bibfnamefont{Y.}~\bibnamefont{Qi}},
  \bibinfo{author}{\bibfnamefont{Z.~Y.} \bibnamefont{Meng}}, \bibnamefont{and}
  \bibinfo{author}{\bibfnamefont{L.}~\bibnamefont{Fu}},
  \bibinfo{journal}{Physical Review B} \textbf{\bibinfo{volume}{95}}
  (\bibinfo{year}{2017}), ISSN \bibinfo{issn}{24699969}.

\bibitem[{\citenamefont{{Liu} et~al.}(2016)\citenamefont{{Liu}, {Shen}, {Qi},
  {Meng}, and {Fu}}}]{Liu2016SLMCFermion}
\bibinfo{author}{\bibfnamefont{J.}~\bibnamefont{{Liu}}},
  \bibinfo{author}{\bibfnamefont{H.}~\bibnamefont{{Shen}}},
  \bibinfo{author}{\bibfnamefont{Y.}~\bibnamefont{{Qi}}},
  \bibinfo{author}{\bibfnamefont{Z.~Y.} \bibnamefont{{Meng}}},
  \bibnamefont{and} \bibinfo{author}{\bibfnamefont{L.}~\bibnamefont{{Fu}}},
  \bibinfo{journal}{arXiv e-prints} \bibinfo{eid}{arXiv:1611.09364}
  (\bibinfo{year}{2016}), \eprint{1611.09364}.

\bibitem[{\citenamefont{Nagai et~al.}(2017)\citenamefont{Nagai, Shen, Qi, Liu,
  and Fu}}]{Nagai2017SLMCCT}
\bibinfo{author}{\bibfnamefont{Y.}~\bibnamefont{Nagai}},
  \bibinfo{author}{\bibfnamefont{H.}~\bibnamefont{Shen}},
  \bibinfo{author}{\bibfnamefont{Y.}~\bibnamefont{Qi}},
  \bibinfo{author}{\bibfnamefont{J.}~\bibnamefont{Liu}}, \bibnamefont{and}
  \bibinfo{author}{\bibfnamefont{L.}~\bibnamefont{Fu}},
  \bibinfo{journal}{Physical Review B} \textbf{\bibinfo{volume}{96}},
  \bibinfo{pages}{1} (\bibinfo{year}{2017}), ISSN \bibinfo{issn}{24699969}.

\bibitem[{\citenamefont{Shen et~al.}(2018)\citenamefont{Shen, Liu, and
  Fu}}]{Shen2018SLMCDNN}
\bibinfo{author}{\bibfnamefont{H.}~\bibnamefont{Shen}},
  \bibinfo{author}{\bibfnamefont{J.}~\bibnamefont{Liu}}, \bibnamefont{and}
  \bibinfo{author}{\bibfnamefont{L.}~\bibnamefont{Fu}},
  \bibinfo{journal}{Physical Review B} \textbf{\bibinfo{volume}{97}}
  (\bibinfo{year}{2018}), ISSN \bibinfo{issn}{24699969}.

\bibitem[{\citenamefont{Li and Wang}(2018)}]{Li2018}
\bibinfo{author}{\bibfnamefont{S.~H.} \bibnamefont{Li}} \bibnamefont{and}
  \bibinfo{author}{\bibfnamefont{L.}~\bibnamefont{Wang}},
  \bibinfo{journal}{Physical Review Letters} \textbf{\bibinfo{volume}{121}},
  \bibinfo{pages}{1} (\bibinfo{year}{2018}), ISSN \bibinfo{issn}{10797114}.

\bibitem[{\citenamefont{{No{\'e}} and {Wu}}(2018)}]{NoeBoltzmannGenerators2018}
\bibinfo{author}{\bibfnamefont{F.}~\bibnamefont{{No{\'e}}}} \bibnamefont{and}
  \bibinfo{author}{\bibfnamefont{H.}~\bibnamefont{{Wu}}},
  \bibinfo{journal}{arXiv e-prints} \bibinfo{eid}{arXiv:1812.01729}
  (\bibinfo{year}{2018}), \eprint{1812.01729}.

\bibitem[{\citenamefont{Song et~al.}(2017)\citenamefont{Song, Zhao, and
  Ermon}}]{Song2017}
\bibinfo{author}{\bibfnamefont{J.}~\bibnamefont{Song}},
  \bibinfo{author}{\bibfnamefont{S.}~\bibnamefont{Zhao}}, \bibnamefont{and}
  \bibinfo{author}{\bibfnamefont{S.}~\bibnamefont{Ermon}}, in
  \emph{\bibinfo{booktitle}{Proceedings in Neural Information Processing
  Systems 2017}} (\bibinfo{year}{2017}), \bibinfo{number}{NIPS}, ISBN
  \bibinfo{isbn}{0008-5472 (Print){\textbackslash}r0008-5472 (Linking)}, ISSN
  \bibinfo{issn}{13000527}, \urlprefix\url{http://arxiv.org/abs/1706.07561}.

\bibitem[{\citenamefont{{Levy} et~al.}(2017)\citenamefont{{Levy}, {Hoffman},
  and {Sohl-Dickstein}}}]{Levy2017GenHMC}
\bibinfo{author}{\bibfnamefont{D.}~\bibnamefont{{Levy}}},
  \bibinfo{author}{\bibfnamefont{M.~D.} \bibnamefont{{Hoffman}}},
  \bibnamefont{and}
  \bibinfo{author}{\bibfnamefont{J.}~\bibnamefont{{Sohl-Dickstein}}},
  \bibinfo{journal}{arXiv e-prints} \bibinfo{eid}{arXiv:1711.09268}
  (\bibinfo{year}{2017}), \eprint{1711.09268}.

\bibitem[{\citenamefont{Urban and Pawlowski}(2018)}]{Urban2018}
\bibinfo{author}{\bibfnamefont{J.~M.} \bibnamefont{Urban}} \bibnamefont{and}
  \bibinfo{author}{\bibfnamefont{J.~M.} \bibnamefont{Pawlowski}}
  (\bibinfo{year}{2018}), \eprint{1811.03533}.

\bibitem[{\citenamefont{Zhou et~al.}(2019)\citenamefont{Zhou, Endrődi, Pang,
  and Stöcker}}]{Zhou2018}
\bibinfo{author}{\bibfnamefont{K.}~\bibnamefont{Zhou}},
  \bibinfo{author}{\bibfnamefont{G.}~\bibnamefont{Endrődi}},
  \bibinfo{author}{\bibfnamefont{L.-G.} \bibnamefont{Pang}}, \bibnamefont{and}
  \bibinfo{author}{\bibfnamefont{H.}~\bibnamefont{Stöcker}},
  \bibinfo{journal}{Phys. Rev.} \textbf{\bibinfo{volume}{D100}},
  \bibinfo{pages}{011501} (\bibinfo{year}{2019}), \eprint{1810.12879}.

\bibitem[{\citenamefont{Tierney}(1994)}]{tierney1994}
\bibinfo{author}{\bibfnamefont{L.}~\bibnamefont{Tierney}},
  \bibinfo{journal}{Ann. Statist.} \textbf{\bibinfo{volume}{22}},
  \bibinfo{pages}{1701} (\bibinfo{year}{1994}),
  \urlprefix\url{https://doi.org/10.1214/aos/1176325750}.

\bibitem[{\citenamefont{Hastings}(1970)}]{hastings70}
\bibinfo{author}{\bibfnamefont{W.~K.} \bibnamefont{Hastings}},
  \bibinfo{journal}{Biometrika} \textbf{\bibinfo{volume}{57}},
  \bibinfo{pages}{97} (\bibinfo{year}{1970}), ISSN \bibinfo{issn}{00063444},
  \urlprefix\url{http://www.jstor.org/stable/2334940}.

\bibitem[{\citenamefont{Rezende and Mohamed}(2015)}]{rezende2015nf}
\bibinfo{author}{\bibfnamefont{D.~J.} \bibnamefont{Rezende}} \bibnamefont{and}
  \bibinfo{author}{\bibfnamefont{S.}~\bibnamefont{Mohamed}},
  \bibinfo{journal}{arXiv preprint arXiv:1505.05770}  (\bibinfo{year}{2015}).

\bibitem[{\citenamefont{Dinh et~al.}(2016)\citenamefont{Dinh, Sohl-Dickstein,
  and Bengio}}]{Dinh2016}
\bibinfo{author}{\bibfnamefont{L.}~\bibnamefont{Dinh}},
  \bibinfo{author}{\bibfnamefont{J.}~\bibnamefont{Sohl-Dickstein}},
  \bibnamefont{and} \bibinfo{author}{\bibfnamefont{S.}~\bibnamefont{Bengio}}
  (\bibinfo{year}{2016}), \urlprefix\url{http://arxiv.org/abs/1605.08803}.

\bibitem[{\citenamefont{Oord et~al.}(2017)\citenamefont{Oord, Li, Babuschkin,
  Simonyan, Vinyals, Kavukcuoglu, Driessche, Lockhart, Cobo, Stimberg
  et~al.}}]{Oord2017WaveNet}
\bibinfo{author}{\bibfnamefont{A.~v.~d.} \bibnamefont{Oord}},
  \bibinfo{author}{\bibfnamefont{Y.}~\bibnamefont{Li}},
  \bibinfo{author}{\bibfnamefont{I.}~\bibnamefont{Babuschkin}},
  \bibinfo{author}{\bibfnamefont{K.}~\bibnamefont{Simonyan}},
  \bibinfo{author}{\bibfnamefont{O.}~\bibnamefont{Vinyals}},
  \bibinfo{author}{\bibfnamefont{K.}~\bibnamefont{Kavukcuoglu}},
  \bibinfo{author}{\bibfnamefont{G.~v.~d.} \bibnamefont{Driessche}},
  \bibinfo{author}{\bibfnamefont{E.}~\bibnamefont{Lockhart}},
  \bibinfo{author}{\bibfnamefont{L.~C.} \bibnamefont{Cobo}},
  \bibinfo{author}{\bibfnamefont{F.}~\bibnamefont{Stimberg}},
  \bibnamefont{et~al.} (\bibinfo{year}{2017}),
  \urlprefix\url{http://arxiv.org/abs/1711.10433}.

\bibitem[{\citenamefont{Zhang et~al.}(2018)\citenamefont{Zhang, E, and
  Wang}}]{DBLP:journals/corr/abs-1809-10188}
\bibinfo{author}{\bibfnamefont{L.}~\bibnamefont{Zhang}},
  \bibinfo{author}{\bibfnamefont{W.}~\bibnamefont{E}}, \bibnamefont{and}
  \bibinfo{author}{\bibfnamefont{L.}~\bibnamefont{Wang}},
  \bibinfo{journal}{CoRR} \textbf{\bibinfo{volume}{abs/1809.10188}}
  (\bibinfo{year}{2018}), \eprint{1809.10188},
  \urlprefix\url{http://arxiv.org/abs/1809.10188}.

\bibitem[{\citenamefont{Kingma and Ba}(2015)}]{Kingma2015Adam}
\bibinfo{author}{\bibfnamefont{D.~P.} \bibnamefont{Kingma}} \bibnamefont{and}
  \bibinfo{author}{\bibfnamefont{J.~L.} \bibnamefont{Ba}}, in
  \emph{\bibinfo{booktitle}{International Conference on Learning
  Representations (ICLR)}} (\bibinfo{year}{2015}).

\bibitem[{\citenamefont{Nesterov}(1983)}]{nesterov1983}
\bibinfo{author}{\bibfnamefont{Y.~E.} \bibnamefont{Nesterov}},
  \bibinfo{journal}{Dokl. Akad. Nauk SSSR} \textbf{\bibinfo{volume}{269}},
  \bibinfo{pages}{543} (\bibinfo{year}{1983}),
  \urlprefix\url{https://ci.nii.ac.jp/naid/10029946121/en/}.

\bibitem[{\citenamefont{Wolff}(2004)}]{Wolff2003}
\bibinfo{author}{\bibfnamefont{U.}~\bibnamefont{Wolff}}
  (\bibinfo{collaboration}{ALPHA}), \bibinfo{journal}{Comput. Phys. Commun.}
  \textbf{\bibinfo{volume}{156}}, \bibinfo{pages}{143} (\bibinfo{year}{2004}),
  \bibinfo{note}{[Erratum: Comput. Phys. Commun.176,383(2007)]},
  \eprint{hep-lat/0306017}.

\bibitem[{\citenamefont{Vierhaus}(2010)}]{vierhaus}
\bibinfo{author}{\bibfnamefont{I.}~\bibnamefont{Vierhaus}}, Ph.D. thesis,
  \bibinfo{school}{Humboldt University of Berlin} (\bibinfo{year}{2010}).

\bibitem[{\citenamefont{Nair and Hinton}(2010)}]{Nair2010}
\bibinfo{author}{\bibfnamefont{V.}~\bibnamefont{Nair}} \bibnamefont{and}
  \bibinfo{author}{\bibfnamefont{G.~E.} \bibnamefont{Hinton}}, in
  \emph{\bibinfo{booktitle}{ICML}} (\bibinfo{year}{2010}).

\bibitem[{\citenamefont{Wolff}(1996)}]{Wolff1996}
\bibinfo{author}{\bibfnamefont{U.}~\bibnamefont{Wolff}},
  \emph{\bibinfo{title}{Numerical Simulation in Quantum Field Theory}}
  (\bibinfo{publisher}{Springer Berlin Heidelberg}, \bibinfo{address}{Berlin,
  Heidelberg}, \bibinfo{year}{1996}), pp. \bibinfo{pages}{245--257}, ISBN
  \bibinfo{isbn}{978-3-642-85238-1},
  \urlprefix\url{https://doi.org/10.1007/978-3-642-85238-1_14}.

\bibitem[{\citenamefont{Duane et~al.}(1987)\citenamefont{Duane, Kennedy,
  Pendleton, and Roweth}}]{Duane1987HMC}
\bibinfo{author}{\bibfnamefont{S.}~\bibnamefont{Duane}},
  \bibinfo{author}{\bibfnamefont{A.}~\bibnamefont{Kennedy}},
  \bibinfo{author}{\bibfnamefont{B.~J.} \bibnamefont{Pendleton}},
  \bibnamefont{and} \bibinfo{author}{\bibfnamefont{D.}~\bibnamefont{Roweth}},
  \bibinfo{journal}{Physics Letters B} \textbf{\bibinfo{volume}{195}},
  \bibinfo{pages}{216 } (\bibinfo{year}{1987}), ISSN \bibinfo{issn}{0370-2693},
  \urlprefix\url{http://www.sciencedirect.com/science/article/pii/037026938791197X}.

\bibitem[{\citenamefont{Goodman and Sokal}(1989)}]{GoodmanSokal89}
\bibinfo{author}{\bibfnamefont{J.}~\bibnamefont{Goodman}} \bibnamefont{and}
  \bibinfo{author}{\bibfnamefont{A.~D.} \bibnamefont{Sokal}},
  \bibinfo{journal}{Phys. Rev. D} \textbf{\bibinfo{volume}{40}},
  \bibinfo{pages}{2035} (\bibinfo{year}{1989}),
  \urlprefix\url{https://link.aps.org/doi/10.1103/PhysRevD.40.2035}.

\bibitem[{\citenamefont{Batrouni and Svetitsky}(1987)}]{Batrouni87}
\bibinfo{author}{\bibfnamefont{G.~G.} \bibnamefont{Batrouni}} \bibnamefont{and}
  \bibinfo{author}{\bibfnamefont{B.}~\bibnamefont{Svetitsky}},
  \bibinfo{journal}{Phys. Rev. B} \textbf{\bibinfo{volume}{36}},
  \bibinfo{pages}{5647} (\bibinfo{year}{1987}),
  \urlprefix\url{https://link.aps.org/doi/10.1103/PhysRevB.36.5647}.

\bibitem[{\citenamefont{Brower and Tamayo}(1989)}]{BrowerTamayo89}
\bibinfo{author}{\bibfnamefont{R.~C.} \bibnamefont{Brower}} \bibnamefont{and}
  \bibinfo{author}{\bibfnamefont{P.}~\bibnamefont{Tamayo}},
  \bibinfo{journal}{Phys. Rev. Lett.} \textbf{\bibinfo{volume}{62}},
  \bibinfo{pages}{1087} (\bibinfo{year}{1989}),
  \urlprefix\url{https://link.aps.org/doi/10.1103/PhysRevLett.62.1087}.

\bibitem[{\citenamefont{Dinh et~al.}(2014)\citenamefont{Dinh, Krueger, and
  Bengio}}]{Dinh2015}
\bibinfo{author}{\bibfnamefont{L.}~\bibnamefont{Dinh}},
  \bibinfo{author}{\bibfnamefont{D.}~\bibnamefont{Krueger}}, \bibnamefont{and}
  \bibinfo{author}{\bibfnamefont{Y.}~\bibnamefont{Bengio}},
  \textbf{\bibinfo{volume}{1}}, \bibinfo{pages}{1} (\bibinfo{year}{2014}), ISSN
  \bibinfo{issn}{1410.8516}, \urlprefix\url{http://arxiv.org/abs/1410.8516}.

\bibitem[{\citenamefont{Kingma and Dhariwal}(2018)}]{Kingma2018}
\bibinfo{author}{\bibfnamefont{D.~P.} \bibnamefont{Kingma}} \bibnamefont{and}
  \bibinfo{author}{\bibfnamefont{P.}~\bibnamefont{Dhariwal}}, in
  \emph{\bibinfo{booktitle}{Neural Information Processing Systems}}
  (\bibinfo{year}{2018}), pp. \bibinfo{pages}{1--15}, ISBN
  \bibinfo{isbn}{0769501850}, ISSN \bibinfo{issn}{10059040},
  \urlprefix\url{http://arxiv.org/abs/1807.03039}.

\bibitem[{\citenamefont{Hoogeboom et~al.}(2019)\citenamefont{Hoogeboom, Berg,
  and Welling}}]{Hoogeboom2019}
\bibinfo{author}{\bibfnamefont{E.}~\bibnamefont{Hoogeboom}},
  \bibinfo{author}{\bibfnamefont{R.~V.~D.} \bibnamefont{Berg}},
  \bibnamefont{and} \bibinfo{author}{\bibfnamefont{M.}~\bibnamefont{Welling}}
  (\bibinfo{year}{2019}), \urlprefix\url{http://arxiv.org/abs/1901.11137}.

\bibitem[{\citenamefont{Grathwohl et~al.}(2019)\citenamefont{Grathwohl, Chen,
  Bettencourt, Sutskever, and Duvenaud}}]{Grathwohl2019}
\bibinfo{author}{\bibfnamefont{W.}~\bibnamefont{Grathwohl}},
  \bibinfo{author}{\bibfnamefont{R.~T.~Q.} \bibnamefont{Chen}},
  \bibinfo{author}{\bibfnamefont{J.}~\bibnamefont{Bettencourt}},
  \bibinfo{author}{\bibfnamefont{I.}~\bibnamefont{Sutskever}},
  \bibnamefont{and} \bibinfo{author}{\bibfnamefont{D.}~\bibnamefont{Duvenaud}},
  in \emph{\bibinfo{booktitle}{ICLR}} (\bibinfo{year}{2019}), pp.
  \bibinfo{pages}{1--13}, ISBN \bibinfo{isbn}{9783901882760}, ISSN
  \bibinfo{issn}{16879139}, \urlprefix\url{http://arxiv.org/abs/1810.01367}.

\bibitem[{\citenamefont{Li and Grathwohl}(2018)}]{Grathwohl2018}
\bibinfo{author}{\bibfnamefont{X.}~\bibnamefont{Li}} \bibnamefont{and}
  \bibinfo{author}{\bibfnamefont{W.}~\bibnamefont{Grathwohl}}, pp.
  \bibinfo{pages}{1--7} (\bibinfo{year}{2018}),
  \urlprefix\url{http://bayesiandeeplearning.org/2018/papers/37.pdf}.

\bibitem[{\citenamefont{{Cohen} et~al.}(2019)\citenamefont{{Cohen}, {Weiler},
  {Kicanaoglu}, and {Welling}}}]{Cohen2019GaugeEquiv}
\bibinfo{author}{\bibfnamefont{T.~S.} \bibnamefont{{Cohen}}},
  \bibinfo{author}{\bibfnamefont{M.}~\bibnamefont{{Weiler}}},
  \bibinfo{author}{\bibfnamefont{B.}~\bibnamefont{{Kicanaoglu}}},
  \bibnamefont{and}
  \bibinfo{author}{\bibfnamefont{M.}~\bibnamefont{{Welling}}},
  \bibinfo{journal}{arXiv e-prints} \bibinfo{eid}{arXiv:1902.04615}
  (\bibinfo{year}{2019}), \eprint{1902.04615}.

\bibitem[{\citenamefont{Valleau and Card}(1972)}]{ValleauFreeEnergy72}
\bibinfo{author}{\bibfnamefont{J.~P.} \bibnamefont{Valleau}} \bibnamefont{and}
  \bibinfo{author}{\bibfnamefont{D.~N.} \bibnamefont{Card}},
  \bibinfo{journal}{The Journal of Chemical Physics}
  \textbf{\bibinfo{volume}{57}}, \bibinfo{pages}{5457} (\bibinfo{year}{1972}),
  \eprint{https://doi.org/10.1063/1.1678245},
  \urlprefix\url{https://doi.org/10.1063/1.1678245}.

\end{thebibliography}

